% =============================================================================
\section{Security Analysis}
\label{sec:security}
% =============================================================================

This section identifies and addresses the principal security considerations for
the PURECORTEX protocol, covering arithmetic precision, reentrancy, LP lock
immutability, oracle-free design, contract sovereignty, and governance safety.

\subsection{Integer Precision in Bonding Curve Arithmetic}
\label{subsec:precision}

The bonding curve cost function (Equation~\eqref{eq:cost-closed}) involves
multiplication and division operations that are susceptible to truncation
errors when implemented in integer arithmetic (Algorand's AVM does not support
floating-point operations).

\subsubsection{Truncation Bug}

A naive implementation of the cost function computes:
\begin{equation}
  C(s, n) = n \cdot B + \frac{S \cdot (2sn + n^2)}{2}
  \label{eq:naive-cost}
\end{equation}

If $S$ is expressed in microCORTEX per token, the intermediate product
$S \cdot (2sn + n^2)$ can overflow a 64-bit unsigned integer for large values
of $s$ and $n$. Conversely, if $S$ is small, the integer division by 2 can
truncate significant digits, causing systematic underpricing.

\subsubsection{Fix: Scaled Arithmetic}

The implementation uses a scaling factor $\Lambda = 10^{12}$ to maintain
precision:

\begin{equation}
  C_{\text{scaled}}(s, n) = \frac{n \cdot B \cdot \Lambda + S \cdot (2sn + n^2) \cdot \Lambda / 2}{\Lambda}
  \label{eq:scaled-cost}
\end{equation}

All intermediate computations are performed in $\Lambda$-scaled space, with
a single final division to return to the native unit. The scaling factor is
chosen such that:
\begin{enumerate}
  \item No intermediate product exceeds $2^{128}$ (Algorand's AVM supports
    128-bit intermediate arithmetic via the \texttt{wide\_ratio} opcode).
  \item The truncation error is bounded by $1 / \Lambda = 10^{-12}$ CORTEX,
    which is negligible relative to the token's 6-decimal precision.
\end{enumerate}

\subsection{Reentrancy Elimination via Atomic Transactions}
\label{subsec:reentrancy}

Reentrancy attacks---wherein a malicious contract re-enters a vulnerable
function during execution---are a well-documented class of vulnerabilities in
EVM-based smart contracts. The PURECORTEX protocol is architecturally immune
to reentrancy due to Algorand's execution model:

\begin{enumerate}
  \item \textbf{Atomic transaction groups.} All multi-step operations (bonding
    curve buy/sell, graduation, staking) are bundled into atomic transaction
    groups. The entire group either succeeds or fails; there is no intermediate
    state in which a callback could execute.
  \item \textbf{No external calls during execution.} Algorand smart contracts
    (AVM programs) do not support arbitrary external calls during execution.
    Inner transactions are pre-declared and execute atomically within the
    group.
  \item \textbf{No fallback functions.} Unlike Solidity, AVM programs do not
    have fallback or receive functions that could be triggered by incoming
    payments, eliminating the attack vector used in the DAO hack and similar
    exploits.
\end{enumerate}

\subsection{LogicSig LP Lock Immutability}
\label{subsec:logicsig}

The LP tokens generated during graduation are locked in a LogicSig account
with the following security properties:

\begin{enumerate}
  \item \textbf{Statelessness.} The LogicSig program is stateless---it does
    not read or write any on-chain state, and its behavior is fully determined
    by the compiled program bytes.
  \item \textbf{Immutability.} Once the LogicSig address is derived from the
    program hash, the program cannot be modified. Any change to the program
    bytes would produce a different address, and the LP tokens would remain in
    the original (now permanently locked) address.
  \item \textbf{No admin keys.} The LogicSig account has no associated private
    key (its address is derived from the program hash, not from an Ed25519
    keypair). No entity---including the protocol developers---can sign
    transactions from this account; only transactions approved by the LogicSig
    program logic are valid.
  \item \textbf{Time enforcement.} The program checks
    $\texttt{global.LatestTimestamp} \geq t_{\text{graduation}} + 10 \cdot
    365.25 \cdot 24 \cdot 3600$ before approving any withdrawal. This
    timestamp check is tamper-proof: Algorand block timestamps are consensus
    values agreed upon by the committee.
\end{enumerate}

\subsection{Oracle-Free Design}
\label{subsec:oracle-free}

The PURECORTEX protocol does not depend on external price oracles for any core
mechanism:

\begin{table}[H]
\centering
\caption{Oracle-free design: all price inputs are derived from on-chain data.}
\label{tab:oracle-free}
\begin{tabular}{@{}lll@{}}
\toprule
\textbf{Mechanism} & \textbf{Price Input} & \textbf{Source} \\
\midrule
Bonding Curve & $P(q) = B + Sq$ & Deterministic formula \\
Dynamic Fees & $\sigma$ (realized vol) & On-chain trade history \\
Buyback & CORTEX market price & Tinyman pool reserves \\
\veCORTEX{} & Lock amount in CORTEX & On-chain state \\
Composability & x402 payment amount & Transaction payload \\
\bottomrule
\end{tabular}
\end{table}

This eliminates the following attack vectors:
\begin{itemize}
  \item \textbf{Oracle manipulation.} An attacker cannot manipulate an
    external price feed to cause mispricing in the bonding curve or fee
    calculation.
  \item \textbf{Oracle downtime.} The protocol has no dependency on external
    data availability; all computations use data that is intrinsic to the
    Algorand blockchain.
  \item \textbf{Flash loan oracle attacks.} Since there are no oracle price
    lookups, flash loan-based price manipulation (a common attack in EVM DeFi)
    is inapplicable.
\end{itemize}

\subsection{Contract Sovereignty}
\label{subsec:sovereignty}

The PURECORTEX smart contracts are designed to be self-sovereign after deployment:

\begin{enumerate}
  \item \textbf{No admin keys.} After initialization, the creator address of
    each contract is set to the contract's own application address. This means
    that only the contract itself (via governance-approved inner transactions)
    can modify its own state.
  \item \textbf{Governance-gated upgrades.} Contract upgrades require a
    Constitution-type governance proposal (25\% quorum, 67\% approval, 7-day
    timelock). The upgrade mechanism uses Algorand's
    \texttt{update\_application} inner transaction, which is only executable
    from within the contract itself.
  \item \textbf{Immutable core logic.} Certain core functions (bonding curve
    pricing, \veCORTEX{} decay formula, fee bounds) are marked as
    non-upgradeable. Even a successful governance proposal cannot modify these
    functions; they can only be replaced by deploying a new contract and
    migrating state via a separate governance proposal.
\end{enumerate}

\subsection{Tri-Brain Governance Safety}
\label{subsec:tri-brain}

The Senator's tri-brain architecture provides defense in depth for governance
proposals:

\begin{enumerate}
  \item \textbf{Analytical check.} The analytical brain verifies that the
    proposed parameter change increases the welfare function $W$
    (Equation~\eqref{eq:welfare}) by at least the significance threshold.
  \item \textbf{Constitutional check.} The constitutional brain evaluates the
    proposal against all seven articles of the constitution
    (Appendix~\ref{app:constitution}), rejecting proposals that violate
    constitutional constraints.
  \item \textbf{Consensus requirement.} At least two of three brains must approve before the
    proposal is submitted (2-of-3 majority agreement). A failure to reach majority
    (e.g., the analytical brain identifies a welfare-increasing change that the constitutional brain flags
    as a rights violation while the third brain abstains) results in proposal rejection.
\end{enumerate}

This tri-brain check mechanism prevents two categories of harmful proposals:
\begin{itemize}
  \item Proposals that are welfare-optimal but unconstitutional (e.g.,
    confiscating a minority's staked tokens to redistribute to the majority).
  \item Proposals that are constitutional but welfare-reducing (e.g., a
    technically valid parameter change that would destabilize the fee market).
\end{itemize}

\subsection{Attack Surface Summary}
\label{subsec:attack-summary}

Table~\ref{tab:attack-summary} summarizes the attack surface and mitigations.

\begin{table}[H]
\centering
\caption{Attack surface and mitigations.}
\label{tab:attack-summary}
\begin{tabular}{@{}p{3.5cm}p{4cm}p{4.5cm}@{}}
\toprule
\textbf{Attack Vector} & \textbf{Vulnerability} & \textbf{Mitigation} \\
\midrule
Arithmetic overflow & Bonding curve computation & Scaled arithmetic + wide\_ratio \\
Reentrancy & Multi-step operations & Atomic transaction groups \\
LP rug pull & LP token withdrawal & LogicSig immutable lock \\
Oracle manipulation & Price feed dependency & Oracle-free design \\
Admin key compromise & Privileged access & Self-sovereign contracts \\
Governance capture & Majority control & Constitutional checks + supermajority \\
Sybil composability & Inflated CS scores & $\sqrt{\cdot}$ diminishing returns \\
\bottomrule
\end{tabular}
\end{table}
